\chapter{Classification of results}
\label{ch:classification}

Every named result in the current theory carries a \emph{classification}.
It answers two independent questions, and keeping them apart is the whole
point of this chapter:

\begin{itemize}
\item \textbf{What justifies the statement?} --- is it a primitive
  assumed as a base, a definition true by construction, a theorem with a
  machine proof, or a result trusted on a stated warrant without proof?
\item \textbf{May the proof engine activate it unbidden?} --- whether the
  prover is allowed to apply the result as a rewrite rule during a proof.
\end{itemize}

\noindent
These two questions are \emph{orthogonal}.  The first sorts results into
the four kinds below; the second is a separate matter, discussed in
Section~\ref{sec:firing}.  A common confusion --- ``is this proved?'' ---
conflates them, and most of what follows is aimed at pulling them apart.

%% =========================================================
\section{The four kinds}
\label{sec:four-kinds}

\begin{itemize}
\item \textbf{Axiom.} A primitive of VNB --- the base we build on.  The
  axiom set is small and essentially fixed; see
  Section~\ref{sec:primitives}.  The only reason to add one is the
  discovery that the existing primitives do not capture the full power of
  VNB.  An axiom is assumed, not proved.

\item \textbf{Definition.} True \emph{by construction}: a defining
  equation, or an unfolding of one.  A definition is \emph{not a theorem}
  --- there is nothing to prove, since it holds by the meaning of the
  symbol it introduces.  Consequently it is never submitted to the VNB
  test and is \emph{not} a member of the Proof Support Set.  It does,
  however, still \textbf{fire as a rewrite} (Section~\ref{sec:firing}):
  the defining equation is exactly the rule that unfolds the symbol.

\item \textbf{Theorem.} Carries a machine proof.  Its proof script runs
  end-to-end and is re-certified by the VNB test
  (Section~\ref{sec:vnb-test}).  This is the only kind whose truth rests
  on a checked derivation inside the system.

\item \textbf{Proof Support Set (PSS).} A theorem \emph{trusted on a
  warrant, without a machine proof} --- and so \emph{excused} from the VNB
  test.  Every PSS member carries credentials: a plain statement of what
  it says plus a short justification of why it is accepted
  (Section~\ref{sec:warrants}).  The PSS is the curated foundational
  engine of obvious-but-tedious facts (finite-sum identities, ordinal and
  set commonplaces) that are not in doubt and not worth the cost of
  mechanising.  A second motive is reuse: a proof typically calls on many
  small, well-established intermediate facts that recur across unrelated
  developments, and collecting them in the PSS makes each available
  everywhere at once rather than re-derived context after context.
\end{itemize}

\begin{remark}
The PSS exists for a practical reason.  Pushing into new
mathematical territory, a proof constantly bumps into basic --- and
genuinely important --- facts whose proofs are well known but tedious and
time-consuming to mechanise.  Stopping to grind each one out would stall
the advance.  The PSS is the mechanism for brushing them aside: assert the
fact on a warrant, keep moving.  That is exactly why it was introduced.
\end{remark}

\begin{remark}
``Asserted'' is not a fifth kind.  A result installed without a proof is
debt against the fixed axiom base: it is destined either to be discharged
(promoted to a \emph{theorem} with a proof) or recognised as a genuine
primitive gap (an \emph{axiom}), or --- if it is an obvious support fact
--- credentialed into the \emph{PSS}.  An asserted result carrying no
warrant at all is a backlog item, not a stable category.
\end{remark}

%% =========================================================
\section{Activation versus justification}
\label{sec:firing}

Whether the prover may apply a result as a rewrite rule is
\emph{independent} of the four-way classification above.  Axioms,
definitions, theorems, and PSS members can \emph{all} be rewrite rules,
and any of them can equally be inert (cited only by name).  Activation is
the business of the \emph{macete} mechanism, a separate registry described
in Section~\ref{sec:macetes}: installing a result as a macete with
\texttt{install-theorem!} is what lets the engine fire it, and that is a
decision about \emph{use}, not about what justifies the statement.

The clearest case is the definition: it is not a PSS member, yet its
defining equation is precisely the rewrite that unfolds the defined
symbol.  Conversely, a fully proved theorem may be left inert if it is too
specialised to be worth firing automatically.

\begin{center}
\begin{tabular}{@{}llll@{}}
\toprule
\textbf{Kind} & \textbf{Believed because} & \textbf{In VNB test?} & \textbf{May fire?}\\
\midrule
Axiom      & assumed (primitive)   & no  & yes\\
Definition & true by construction  & no  & yes\\
Theorem    & machine proof         & yes & yes\\
PSS        & a stated warrant      & no (excused) & yes\\
\bottomrule
\end{tabular}
\end{center}

\noindent
The last column is the same ``yes'' in every row: firing rights cut across
the classification entirely.

%% =========================================================
\section{Warrants}
\label{sec:warrants}

A definition needs no warrant: it is true by construction, so there is
nothing to justify.  A PSS member, by contrast, is an asserted theorem ---
something that \emph{could} be false and is not being proved here --- and
so it must carry credentials.  A \emph{warrant} is that justification: a
comprehensible statement of the fact together with at least a couple of
lines saying why we accept it without a machine proof.

A warrant is attached with

\begin{VNBsnip}
(warrant! 'name 'kind "free-text justification")
\end{VNBsnip}

\noindent
where \texttt{kind} is one of the recognised grades of justification (the
set \texttt{*warrant-kinds*}), weakest to strongest:

\begin{itemize}
\item \texttt{hand-wave} --- a heuristic or plausibility argument;
  convincing, not rigorous.
\item \texttt{well-known} --- a standard textbook fact, asserted without
  argument.
\item \texttt{reference} --- cites a specific source, named in the text.
\item \texttt{informal} --- a rigorous paper-proof exists (often sketched
  in the file comment or the text), just not mechanised in VNB.
\item \texttt{proof} --- a machine-checked VNB proof exists.
\end{itemize}

\begin{remark}
A warrant is a promissory note, not a permanent verdict: any weaker kind
may later be discharged into \texttt{proof}.  Nothing logical depends on
it --- it is metadata for review and triage.  A result with no warrant at
all ranks below even \texttt{hand-wave}.
\end{remark}

\begin{remark}
Warrants propagate.  When a biconditional theorem is installed, its
automatically generated \texttt{-rev} companion inherits the same warrant,
so the reverse-direction rewrite is credentialed identically to the
forward one.
\end{remark}

%% =========================================================
\section{Proof debt: bills and trust tiers}
\label{sec:proof-debt}

A machine proof is only as unconditional as the facts it cites.  When a
proof reaches \texttt{QED} the system computes its \emph{bill}: the set of
\emph{asserted} facts it ultimately rests on.  The bill is transitive.  A
cited \emph{proven} lemma contributes its own bill; a \emph{primitive}
axiom or a \emph{definition} contributes nothing (it is trusted base); an
\emph{asserted} fact --- a PSS member, or an as-yet-unwarranted assertion
--- contributes itself.  Each fact in the bill is a \emph{leaf}.

The bill is printed inline at \texttt{QED} and collected in the ledger
\texttt{reference/PROOF-DEBT.md}:

\begin{itemize}
\item \texttt{proven modulo 0} --- the \emph{empty bill}.  The proof rests
  only on the trusted base: the fixed primitive axioms, the kernel
  inference rules, and definitions.  Nothing asserted stands between the
  theorem and the kernel.  This is the strongest report a proof can make,
  and it carries \emph{no} trust tier.
\item \texttt{proven modulo \{a, b, \ldots\}} --- unconditional
  \emph{given} $a, b, \ldots$.  Discharging every leaf (proving it, or
  turning it into a definition) would collapse the bill to $0$.
\end{itemize}

\noindent
A non-empty bill is stamped with a \emph{trust tier}
\texttt{[trust:\ \textit{K}]}, and $K$ is the warrant of the bill's
\emph{weakest} leaf: a single shaky citation caps the whole proof, however
solid the rest.  The tiers are the warrant kinds of
Section~\ref{sec:warrants} with one addition below the weakest ---
\texttt{none}, an asserted leaf carrying no warrant at all --- so, weakest
to strongest:
\[
  \texttt{none} \;<\; \texttt{hand-wave} \;<\; \texttt{well-known}
  \;<\; \texttt{reference} \;<\; \texttt{informal} \;<\; \texttt{proof}.
\]
Thus \texttt{[trust:\ none]} is the \emph{weakest} thing a proof can
report --- some leaf justifies nothing --- and \texttt{modulo 0}, which
prints no tier at all, is the strongest.  The two are opposite ends of the
same scale, and are easily misread for one another; the ledger's header
carries this same legend for exactly that reason.

\begin{remark}
The tier is a moving target by design.  Discharging a \texttt{none} or
\texttt{hand-wave} leaf into a warranted or proven fact lifts the tier of
every proof that cites it.  The ledger's reverse index ranks the asserted
leaves by how many proofs each one caps, so the highest-leverage debts to
retire are read off the top.
\end{remark}

%% =========================================================
\section{Circular dependency}
\label{sec:proof-cycle}

The bill of Section~\ref{sec:proof-debt} is computed from a \emph{citation
graph}: at each \texttt{QED} the system records which named results the finished
proof cited.  That graph carries a soundness obligation of its own, independent
of the trust tiers: \textbf{no proven theorem may depend --- transitively,
through its citations --- on itself.}

The failure it guards against is a circle that no single step betrays.  Assert
$X$; prove $Y$ citing $X$; then ``prove'' $X$ citing $Y$.  Each proof is locally
valid, yet together they rest on nothing --- $X$ is believed because of $Y$, and
$Y$ because of $X$.  A bill, a frozen snapshot taken at \texttt{QED}, cannot
reveal such a loop, so the check walks the \emph{live} citation graph instead.

\texttt{(proof-cycle-check)} returns every dependency cycle among the proven
theorems --- the empty list when there are none.  \texttt{load.scm} runs it as a
load-time gate and, on any surviving cycle, \textbf{aborts the load} rather than
merely warn: a genuine circular proof is a soundness bug and is treated as one.

Two kinds of citation edge are deliberately \emph{not} dependencies, and are
excluded when the graph is recorded.  Without the exclusion an ordinary
well-founded induction would read as a cycle:

\begin{itemize}
\item \textbf{A proof's citation of its own name.}  A proof by induction applies
  its induction hypothesis --- the theorem's own statement at a smaller
  instance --- as a rewrite; likewise a structure predicate is unfolded by a
  macete bearing the very name the finished theorem will carry.  Neither is a
  dependence of the theorem on itself.
\item \textbf{A citation that is a definition when it is recorded.}  Unfolding a
  definition adds no debt (Section~\ref{sec:proof-debt}) and no dependency; a
  definitional macete that happens to share a later theorem's name is excluded
  on the same ground.
\end{itemize}

\noindent
Excluding these leaves every genuine circular dependency intact --- such a loop
runs through \emph{proven} citations, which are always kept --- so the gate
neither cries wolf on sound induction nor lets a real circle through.

%% =========================================================
\section{The VNB test}
\label{sec:vnb-test}

The \emph{VNB test} is the standing examination that certifies the
\textbf{theorem} bucket.  It re-proves, in sequence, every result that has
a runnable proof script, and reports whether each still reaches
\texttt{QED}.  Definitions and PSS members are excused by construction:
the former have nothing to prove, the latter are accepted on a warrant.
Run it from the prover root:

\begin{VNBsnip}
./vnb-test
\end{VNBsnip}

\noindent
It harvests the proofs that run at load time and the standalone proof
scripts, runs each, and writes a \emph{certification stamp}
(\texttt{reference/certification.scm}) recording the date and which results
passed, together with a human-readable report
(\texttt{reference/VNB-TEST.md}).  The next time the prover loads, the
stamp is read back, so a result's status can report ``certified by the VNB
test on \textit{date}'' rather than merely ``has a proof script.''

\begin{remark}
The VNB test is \emph{not} the unit test suite.
\texttt{test-suite.scm} is a fast feature-and-regression suite of several
hundred cases run constantly during development; it checks that the
machinery behaves, not that the library's theorems still derive.  The VNB
test is the slower, periodic exam that re-runs the mathematics
itself --- re-deriving every proof in the library, as any serious proof
system must.  Keep both: they answer different questions.
\end{remark}

%% =========================================================
\section{Inspecting a result}
\label{sec:status-query}

Three commands expose the classification at the REPL.

\begin{itemize}
\item \texttt{(status 'name)} prints, on one screen, a result's
  provenance, whether it is a PSS member, any warrant it carries, and its
  certification state (``certified by the VNB test on \textit{date}'' or
  ``not certified since the last run'').

\item \texttt{(status-audit)} writes \texttt{reference/STATUS-AUDIT.md}: a
  census of the whole library by kind, flagging hygiene issues --- asserted
  results with no warrant, theorem-shaped names that have drifted out of
  the theorem bucket, and the like.

\item \texttt{(pss-letter 'name)} drafts a plain-English ``recommendation
  letter'' for a candidate PSS member: an English rendering of what it
  says, its current standing, and an assessment against the admission
  gates (true; a sound, terminating rewrite shape; general rather than
  bespoke; non-redundant).  It is an aid to deciding whether a result
  belongs in the curated support set, not an automatic admission.
\end{itemize}
